\documentclass[a4paper]{article}

\usepackage{epsfig}

\title{\epsfig{file=coa-hor.eps,width=6cm}\vspace{3cm}\\WEKA Explorer User Guide\\for Version 3-4}
\author{Richard Kirkby\\Eibe Frank}

\setcounter{secnumdepth}{1}

\begin{document}

\begin{titlepage}

\maketitle
\thispagestyle{empty}

\center
\vspace{8cm}

\copyright 2002-2005 University of Waikato

\end{titlepage}

\tableofcontents

\newpage

\section{Launching WEKA}

The WEKA GUI Chooser window is used to launch WEKA's graphical
environments. At the bottom of the window are four buttons:

\begin{enumerate}
\item \textbf{Simple CLI}.  Provides a simple command-line interface
that allows direct execution of WEKA commands for operating systems
that do not provide their own command line interface. 
\item \textbf{Explorer}.  An environment for exploring data with
WEKA. 
\item \textbf{Experimenter}.
An environment for performing experiments and conducting statistical tests
between learning schemes. 
\item \textbf{Knowledge Flow}.  This environment supports essentially
the same functions as the Explorer but with a drag-and-drop
interface. One advantage is that it supports incremental learning.
\end{enumerate}

If you launch WEKA from a terminal window, some text begins scrolling in the
terminal.  Ignore this text unless something goes wrong, in which case it can
help in tracking down the cause.

This User Manual, which is also available online on the \textit{WekaDoc Wiki} 
\cite{wekadoc}, focuses on using the Explorer but does not explain
the individual data preprocessing tools and learning algorithms in
WEKA. For more information on the various filters and learning methods
in WEKA, see the book {\em Data Mining} \cite{witten}.

\section{The WEKA Explorer}

\subsection{Section Tabs}

At the very top of the window, just below the title bar, is a row of
tabs. When the Explorer is first started only the first tab is active;
the others are greyed out. This is because it is necessary to open
(and potentially pre-process) a data set before starting to explore
the data.

The tabs are as follows:

\begin{enumerate}
\item \textbf{Preprocess}.
Choose and modify the data being acted on.
\item \textbf{Classify}.
Train and test learning schemes that classify or perform regression.
\item \textbf{Cluster}.
Learn clusters for the data.
\item \textbf{Associate}.
Learn association rules for the data.
\item \textbf{Select attributes}.
Select the most relevant attributes in the data.
\item \textbf{Visualize}.
View an interactive 2D plot of the data.
\end{enumerate}
\noindent
Once the tabs are active, clicking on them flicks between different
screens, on which the respective actions can be performed.  The bottom
area of the window (including the status box, the log button, and the
Weka bird) stays visible regardless of which section you are in.

\subsection{Status Box}

The status box appears at the very bottom of the window. It displays
messages that keep you informed about what's going on. For example, if
the Explorer is busy loading a file, the status box will say that.

\textbf{TIP}---right-clicking the mouse anywhere inside the status box brings
up a little menu. The menu gives two options: 

\begin{enumerate}
\item \textbf{Memory information}.
Display in the log box the amount of memory available to WEKA.
\item \textbf{Run garbage collector}.
Force the Java garbage collector to search for memory that is no longer needed
and free it up, allowing more memory for new tasks. Note that the garbage
collector is constantly running as a background task anyway.
\end{enumerate}

\subsection{Log Button}

Clicking on this button brings up a separate window containing a scrollable text
field. Each line of text is stamped with the time it was entered into the
log. As you perform actions in WEKA, the log keeps a record of what has
happened.

\subsection{WEKA Status Icon}

To the right of the status box is the WEKA status icon. When no processes are
running, the bird sits down and takes a nap. The number beside the $\times$
symbol gives the number of concurrent processes running.  When the system is
idle it is zero, but it increases as the number of processes increases.
When any process is started, the bird gets up and starts moving around. If
it's standing but stops moving for a long time, it's sick: something has gone
wrong!  In that case you should restart the WEKA explorer.

\section{Preprocessing}

\subsection{Opening files}

The first three buttons at the top of the preprocess section enable
you to load data into WEKA:

\begin{enumerate}
\item \textbf{Open file...}.
Brings up a dialog box allowing you to browse for the data file on the local
filesystem.
\item \textbf{Open URL...}.
Asks for a Uniform Resource Locator address for where the data is stored.
\item \textbf{Open DB...}.  Reads data from a database. (Note that to
make this work you might have to edit the file in
weka/experiment/DatabaseUtils.props.)
\end{enumerate}
\noindent
Using the \textbf{Open file...} button you can read files in a variety
of formats: Weka's ARFF format, CSV format, C4.5 format, or serialized
Instances format. ARFF files typically have a {\em .arff\/}
extension, CSV files a {\em .csv\/} extension, C4.5 files a {\em
.data\/} and {\em .names\/} extension, and serialized Instances
objects a {\em .bsi\/} extension.

\subsection{The Current Relation}

Once some data has been loaded, the Preprocess panel shows a variety
of information.  The \textbf{Current relation} box (the ``current
relation'' is the currently loaded data, which can be interpreted as a
single relational table in database terminology)  has three entries:

\begin{enumerate}
\item \textbf{Relation}.
The name of the relation, as given in the file it was loaded from. Filters
(described below) modify the name of a relation. 
\item \textbf{Instances}.
The number of instances (data points/records) in the data.
\item \textbf{Attributes}.
The number of attributes (features) in the data.
\end{enumerate}

\subsection{Working With Attributes}

Below the \textbf{Current relation} box is a box titled \textbf{Attributes}.
There are three buttons, and beneath them is a list of the attributes in the
current relation. The list has three columns:

\begin{enumerate}
\item \textbf{No.}.
A number that identifies the attribute in the order they are specified in the
data file. 
\item \textbf{Selection tick boxes}.
These allow you select which attributes are present in the relation.
\item \textbf{Name}.
The name of the attribute, as it was declared in the data file.
\end{enumerate}

When you click on different rows in the list of attributes, the fields
change in the box to the right titled \textbf{Selected
attribute}. This box displays the characteristics of the currently
highlighted attribute in the list:

\begin{enumerate}
\item \textbf{Name}.
The name of the attribute, the same as that given in the attribute list.
\item \textbf{Type}.
The type of attribute, most commonly Nominal or Numeric.
\item \textbf{Missing}.
The number (and percentage) of instances in the data for which this attribute
is missing (unspecified).
\item \textbf{Distinct}.
The number of different values that the data contains for this attribute.
\item \textbf{Unique}.
The number (and percentage) of instances in the data having a value for this
attribute that no other instances have.
\end{enumerate}
\noindent
Below these statistics is a list showing more information about the
values stored in this attribute, which differ depending on its type.
If the attribute is nominal, the list consists of each possible value
for the attribute along with the number of instances that have that
value.  If the attribute is numeric, the list gives four statistics
describing the distribution of values in the data---the minimum,
maximum, mean and standard deviation.  And below these statistics
there is a colored histogram, color-coded according to the attribute
chosen as the {\it Class} using the box above the histogram. (This box
will bring up a drop-down list of available selections when clicked.)
Note that only nominal {\it Class} attributes will result in a
color-coding.  Finally, after pressing the \textbf{Visualize All}
button, histograms for all the attributes in the data are shown in a
separate witting.

Returning to the attribute list, to begin with all the tick boxes are unticked.
They can be toggled on/off by clicking on them individually.  The three buttons
above can also be used to change the selection:

\begin{enumerate}
\item \textbf{All}.
All boxes are ticked.
\item \textbf{None}.
All boxes are cleared (unticked).
\item \textbf{Invert}.
Boxes that are ticked become unticked and {\em vice versa\/}.
\end{enumerate}

Once the desired attributes have been selected, they can be removed by
clicking the \textbf{Remove} button below the list of attributes.
Note that this can be undone by clicking the \textbf{Undo} button,
which is located next to the \textbf{Edit} button in the top-right
corner of the Preprocess panel.

\subsection{Working With Filters}

The preprocess section allows filters to be defined that transform the
data in various ways.  The \textbf{Filter} box is used to set up the
filters that are required.  At the left of the \textbf{Filter} box is
a \textbf{Choose} button. By clicking this button it is possible to
select one of the filters in Weka. Once a filter has been selected,
its name and options are shown in the field next to the
\textbf{Choose} button. Clicking on this box brings up a
GenericObjectEditor dialog box.

\subsubsection*{The GenericObjectEditor Dialog Box}

The GenericObjectEditor dialog box lets you configure a filter. The
same kind of dialog box is used to configure other objects, such as
classifiers and clusterers (see below). The fields in the window
reflect the available options.  Clicking on any of these gives an
opportunity to alter the filters settings. For example, the setting
may take a text string, in which case you type the string into the
text field provided.  Or it may give a drop-down box listing several
states to choose from.  Or it may do something else, depending on the
information required. Information on the options is provided in a tool
tip if you let the mouse pointer hover of the corresponding
field. More information on the filter and its options can be obtained
by clicking on the \textbf{More} button in the \textbf{About} panel at
the top of the GenericObjectEditor window.

Some objects display a brief description of what they do in an \textbf{About}
box, along with a \textbf{More} button. Clicking on the \textbf{More} button
brings up a window describing what the different options do.

At the bottom of the GenericObjectEditor dialog are four buttons. The first
two, \textbf{Open...} and \textbf{Save...} allow object configurations to be
stored for future use. The \textbf{Cancel} button backs out without remembering
any changes that have been made.  Once you are happy with the object and
settings you have chosen, click \textbf{OK} to return to the main Explorer
window.

\subsubsection*{Applying Filters}

Once you have selected and configured a filter, you can apply it to
the data by pressing the \textbf{Apply} button at the right end of the
\textbf{Filter} panel in the Preprocess panel. The Preprocess panel
will then show the transformed data. The change can be undone by
pressing the \textbf{Undo} button. You can also use the
\textbf{Edit...} button to modify your data manually in a dataset
editor.  Finally, the \textbf{Save...}  button at the top right of the
Preprocess panel saves the current version of the relation in the same
formats available for loading data, allowing it to be kept for future
use.  \\

\noindent \textbf{Note:} Some of the filters behave differently
depending on whether a class attribute has been set or not (using the
box above the histogram, which will bring up a drop-down list of
possible selections when clicked). In particular, the ``supervised
filters'' require a class attribute to be set, and some of the
``unsupervised attribute filters'' will skip the class attribute if
one is set. Note that it is also possible to set {\em Class} to {\em
None}, in which case no class is set.

\section{Classification}

\subsection{Selecting a Classifier}

\label{sec:classifier}
At the top of the classify section is the \textbf{Classifier}
box. This box has a text field that gives the name of the currently
selected classifier, and its options. Clicking on the text box brings
up a GenericObjectEditor dialog box, just the same as for filters,
that you can use to configure the options of the current classifier.
The \textbf{Choose} button allows you to choose one of the classifiers
that are available in WEKA.

\subsection{Test Options}

The result of applying the chosen classifier will be tested according to the
options that are set by clicking in the \textbf{Test options} box.  There are
four test modes:

\begin{enumerate}
\item \textbf{Use training set}.
The classifier is evaluated on how well it predicts the class of the instances
it was trained on. 
\item \textbf{Supplied test set}.
The classifier is evaluated on how well it predicts the class of a set of
instances loaded from a file. Clicking the \textbf{Set...} button brings up a
dialog allowing you to choose the file to test on.
\item \textbf{Cross-validation}.
The classifier is evaluated by cross-validation, using the number of folds that
are entered in the \textbf{Folds} text field. 
\item \textbf{Percentage split}.
The classifier is evaluated on how well it predicts a certain percentage of the
data which is held out for testing. The amount of data held out depends on the
value entered in the \textbf{\%} field.
\end{enumerate}
\noindent
\textbf{Note:} No matter which evaluation method is used, the model
that is output is always the one build from \textbf{\em all} the training data.
\noindent
Further testing options can be set by clicking on the \textbf{More options...}
button:

\begin{enumerate}
\item \textbf{Output model}.
The classification model on the full training set is output so that it can be
viewed, visualized, etc. This option is selected by default.
\item \textbf{Output per-class stats}.  The precision/recall and
true/false statistics for each class are output. This option is also
selected by default.
\item \textbf{Output entropy evaluation measures}.  Entropy evaluation
measures are included in the output. This option is not selected by
default.
\item \textbf{Output confusion matrix}.
The confusion matrix of the classifier's predictions is included in the output. This option is selected by default.
\item \textbf{Store predictions for visualization}.  The classifier's
predictions are remembered so that they can be visualized. This option
is selected by default.
\item \textbf{Output predictions}. The predictions on the evaluation
data are output.  Note that in the case of a cross-validation the
instance numbers do not correspond to the location in the data!
\item \textbf{Cost-sensitive evaluation}.
The errors is evaluated with respect to a cost matrix. The \textbf{Set...}
button allows you to specify the cost matrix used. 
\item \textbf{Random seed for xval / \% Split}.
This specifies the random seed used when randomizing the data before it is
divided up for evaluation purposes.
\end{enumerate}

\subsection{The Class Attribute}

The classifiers in WEKA are designed to be trained to predict a single `class'
attribute, which is the target for prediction. Some classifiers can only learn
nominal classes; others can only learn numeric classes (regression problems);
still others can learn both.

By default, the class is taken to be the last attribute in the data.  If you
want to train a classifier to predict a different attribute, click on the box
below the \textbf{Test options} box to bring up a drop-down list of attributes
to choose from.

\subsection{Training a Classifier}

Once the classifier, test options and class have all been set, the learning
process is started by clicking on the \textbf{Start} button. While the
classifier is busy being trained, the little bird moves around. You can stop
the training process at any time by clicking on the \textbf{Stop} button.

When training is complete, several things happen. The \textbf{Classifier
output} area to the right of the display is filled with text describing the
results of training and testing. A new entry appears in the \textbf{Result
list} box. We look at the result list below; but first we investigate the text
that has been output.

\subsection{The Classifier Output Text}

The text in the \textbf{Classifier output} area has scroll bars allowing you to
browse the results.  Of course, you can also resize the Explorer window to get
a larger display area.  The output is split into several sections:

\begin{enumerate}
\item \textbf{Run information}.
A list of information giving the learning scheme options, relation name,
instances, attributes and test mode that were involved in the process.
\item \textbf{Classifier model (full training set)}.
A textual representation of the classification model that was produced on the
full training data. 
\item The results of the chosen test mode are broken down thus:
\item \textbf{Summary}.
A list of statistics summarizing how accurately the classifier was able to
predict the true class of the instances under the chosen test mode. 
\item \textbf{Detailed Accuracy By Class}.
A more detailed per-class break down of the classifier's prediction accuracy. 
\item \textbf{Confusion Matrix}.
Shows how many instances have been assigned to each class. Elements show the
number of test examples whose actual class is the row and whose predicted class
is the column.
\end{enumerate}

\subsection{The Result List}

After training several classifiers, the result list will contain several
entries.  Left-clicking the entries flicks back and forth between the various
results that have been generated.  Right-clicking an entry invokes a menu
containing these items:

\begin{enumerate}
\item \textbf{View in main window}.
Shows the output in the main window (just like left-clicking the entry).
\item \textbf{View in separate window}.
Opens a new independent window for viewing the results.
\item \textbf{Save result buffer}.
Brings up a dialog allowing you to save a text file containing the textual
output.
\item \textbf{Load model}.
Loads a pre-trained model object from a binary file.
\item \textbf{Save model}.
Saves a model object to a binary file. Objects are saved in Java `serialized
object' form.
\item \textbf{Re-evaluate model on current test set}.
Takes the model that has been built and tests its performance on the data set
that has been specified with the \textbf{Set..} button under the
\textbf{Supplied test set} option.
\item \textbf{Visualize classifier errors}.
Brings up a visualization window that plots the results of classification.
Correctly classified instances are represented by crosses, whereas incorrectly
classified ones show up as squares.
\item \textbf{Visualize tree} or \textbf{Visualize graph}.  Brings up
a graphical representation of the structure of the classifier model,
if possible (i.e. for decision trees or Bayesian networks). The
graph visualization option only appears if a Bayesian network
classifier has been built. In the tree visualizer, you can bring up a
menu by right-clicking a blank area, pan around by dragging the mouse,
and see the training instances at each node by clicking on
it. CTRL-clicking zooms the view out, while SHIFT-dragging a box zooms
the view in. The graph visualizer should be self-explanatory.
\item \textbf{Visualize margin curve}.
Generates a plot illustrating the prediction margin. The margin is defined as
the difference between the probability predicted for the actual class and the
highest probability predicted for the other classes. For example, boosting
algorithms may achieve better performance on test data by increasing the
margins on the training data. 
\item \textbf{Visualize threshold curve}.
Generates a plot illustrating the tradeoffs in prediction that are obtained by
varying the threshold value between classes. For example, with the default
threshold value of 0.5, the predicted probability of `positive' must be greater
than 0.5 for the instance to be predicted as `positive'. The plot can be used
to visualize the precision/recall tradeoff, for ROC curve analysis (true
positive rate {\em vs} false positive rate), and for other types of curves. 
\item \textbf{Visualize cost curve}.
Generates a plot that gives an explicit representation of the expected cost, as
described by \cite{drummond}.
\end{enumerate}
\noindent
Options are greyed out if they do not apply to the specific set of results.

\section{Clustering}

\subsection{Selecting a Clusterer}

By now you will be familiar with the process of selecting and configuring
objects.  Clicking on the clustering scheme listed in the \textbf{Clusterer}
box at the top of the window brings up a GenericObjectEditor dialog with which
to choose a new clustering scheme.

\subsection{Cluster Modes}

The \textbf{Cluster mode} box is used to choose what to cluster and how to
evaluate the results. The first three options are the same as for
classification: \textbf{Use training set}, \textbf{Supplied test set} and
\textbf{Percentage split} (Section~\ref{sec:classifier})---except that now the
data is assigned to clusters instead of trying to predict a specific class.
The fourth mode, \textbf{Classes to clusters evaluation}, compares how well the
chosen clusters match up with a pre-assigned class in the data.  The drop-down
box below this option selects the class, just as in the \textbf{Classify}
panel.

An additional option in the \textbf{Cluster mode} box, the \textbf{Store
clusters for visualization} tick box, determines whether or not it will be
possible to visualize the clusters once training is complete. When dealing
with datasets that are so large that memory becomes a problem it may be
helpful to disable this option.

\subsection{Ignoring Attributes}

Often, some attributes in the data should be ignored when clustering.  The
\textbf{Ignore attributes} button brings up a small window that allows you to
select which attributes are ignored.  Clicking on an attribute in the window
highlights it, holding down the SHIFT key selects a range of consecutive
attributes, and holding down CTRL toggles individual attributes on and off. To
cancel the selection, back out with the \textbf{Cancel} button. To activate it,
click the \textbf{Select} button. The next time clustering is invoked, the
selected attributes are ignored.

\subsection{Learning Clusters}

The \textbf{Cluster} section, like the \textbf{Classify} section, has
\textbf{Start}/\textbf{Stop} buttons, a result text area and a result
list. These all behave just like their classification counterparts.
Right-clicking an entry in the result list brings up a similar menu,
except that it shows only two visualization options: \textbf{Visualize
cluster assignments} and \textbf{Visualize tree}. The latter is grayed
out when it is not applicable.

\section{Associating}

\subsection{Setting Up}

This panel contains schemes for learning association rules, and the
learners are chosen and configured in the same way as the clusterers,
filters, and classifiers in the other panels.

\subsection{Learning Associations}

Once appropriate parameters for the association rule learner bave been
set, click the \textbf{Start} button.  When complete, right-clicking
on an entry in the result list allows the results to be viewed or
saved.

\section{Selecting Attributes}

\subsection{Searching and Evaluating}

Attribute selection involves searching through all possible combinations of
attributes in the data to find which subset of attributes works best for
prediction.  To do this, two objects must be set up: an attribute evaluator and
a search method.  The evaluator determines what method is used to assign a
worth to each subset of attributes.  The search method determines what style of
search is performed.

\subsection{Options}

The \textbf{Attribute Selection Mode} box has two options:

\begin{enumerate}
\item \textbf{Use full training set}.
The worth of the attribute subset is determined using the full set of training
data. 
\item \textbf{Cross-validation}.
The worth of the attribute subset is determined by a process of
cross-validation.  The \textbf{Fold} and \textbf{Seed} fields set the number of
folds to use and the random seed used when shuffling the data.
\end{enumerate}
\noindent
As with \textbf{Classify} (Section~\ref{sec:classifier}), there is a drop-down
box that can be used to specify which attribute to treat as the class.

\subsection{Performing Selection}

Clicking \textbf{Start} starts running the attribute selection
process.  When it is finished, the results are output into the result
area, and an entry is added to the result list.  Right-clicking on the
result list gives several options. The first three, (\textbf{View in
main window}, \textbf{View in separate window} and \textbf{Save result
buffer}), are the same as for the classify panel.  It is also possible
to \textbf{Visualize reduced data}, or if you have used an attribute
transformer such as PrincipalComponents, \textbf{Visualize transformed
data}.

\section{Visualizing}

WEKA's visualization section allows you to visualize 2D plots of the
current relation.  

\subsection{The scatter plot matrix}

When you select the {\em Visualize} panel, it shows a scatter plot
matrix for all the attributes, color coded according to the currently
selected class. It is possible to change the size of each individual
2D plot and the point size, and to randomly jitter the data (to
uncover obscured points). It also possible to change the attribute
used to color the plots, to select only a subset of attributes for
inclusion in the scatter plot matrix, and to sub sample the data. Note
that changes will only come into effect once the \textbf{Update}
button has been pressed.

\subsection{Selectin an individual 2D scatter plot}
 
When you click on a cell in the scatter plot matrix, this will bring
up a separate window with a visualization of the scatter plot you
selected.  (We described above how to visualize particular results in
a separate window---for example, classifier errors---the same
visualization controls are used here.)

Data points are plotted in the main area of the window.  At the top
are two drop-down list buttons for selecting the axes to plot.  The
one on the left shows which attribute is used for the x-axis; the one
on the right shows which is used for the y-axis.

Beneath the x-axis selector is a drop-down list for choosing the
colour scheme.  This allows you to colour the points based on the
attribute selected.  Below the plot area, a legend describes what
values the colours correspond to. If the values are discrete, you can
modify the colour used for each one by clicking on them and making an
appropriate selection in the window that pops up.

To the right of the plot area is a series of horizontal strips. Each
strip represents an attribute, and the dots within it show the
distribution of values of the attribute.  These values are randomly
scattered vertically to help you see concentrations of points.  You
can choose what axes are used in the main graph by clicking on these
strips.  Left-clicking an attribute strip changes the x-axis to that
attribute, whereas right-clicking changes the y-axis. The `X' and `Y'
written beside the strips shows what the current axes are (`B' is used
for `both X and Y').

Above the attribute strips is a slider labelled \textbf{Jitter}, which
is a random displacement given to all points in the plot.  Dragging it
to the right increases the amount of jitter, which is useful for
spotting concentrations of points. Without jitter, a million instances
at the same point would look no different to just a single lonely
instance.

\subsection{Selecting Instances}

There may be situations where it is helpful to select a subset of the
data using the visualization tool. (A special case of this is the
UserClassifier in the {\em Classify} panel, which lets you build your
own classifier by interactively selecting instances.)

Below the y-axis selector button is a drop-down list button for choosing a
selection method.  A group of data points can be selected in four ways:

\begin{enumerate}
\item \textbf{Select Instance}.
Clicking on an individual data point brings up a window listing its attributes.
If more than one point appears at the same location, more than one set of
attributes is shown.
\item \textbf{Rectangle}.
You can create a rectangle, by dragging, that selects the points inside it.
\item \textbf{Polygon}.
You can build a free-form polygon that selects the points inside it. Left-click
to add vertices to the polygon, right-click to complete it. The polygon will
always be closed off by connecting the first point to the last.
\item \textbf{Polyline}.
You can build a polyline that distinguishes the points on one side from those
on the other. Left-click to add vertices to the polyline, right-click to
finish. The resulting shape is open (as opposed to a polygon, which is
always closed).
\end{enumerate}

Once an area of the plot has been selected using \textbf{Rectangle},
\textbf{Polygon} or \textbf{Polyline}, it turns grey.  At this point, clicking
the \textbf{Submit} button removes all instances from the plot except those
within the grey selection area.  Clicking on the \textbf{Clear} button erases
the selected area without affecting the graph.

Once any points have been removed from the graph, the \textbf{Submit} button
changes to a \textbf{Reset} button.  This button undoes all previous removals
and returns you to the original graph with all points included.  Finally,
clicking the \textbf{Save} button allows you to save the currently visible
instances to a new ARFF file.

\begin{thebibliography}{999}
   \bibitem{drummond} Drummond, C. and Holte, R. (2000) Explicitly representing expected cost: An alternative to ROC representation. \textit{Proceedings of the Sixth ACM SIGKDD International Conference on Knowledge Discovery and Data Mining.}
Publishers, San Mateo, CA.
	\bibitem{witten} Witten, I.H. and Frank, E. (2005) \textit{Data Mining: Practical machine
learning tools and techniques. 2nd edition}  Morgan Kaufmann, San
Francisco.
	\bibitem{wekadoc} \textit{WekaDoc} -- \texttt{http://weka.sourceforge.net/wekadoc/}
\end{thebibliography}

\end{document}
